\section{Results}
\label{sec:results}

\paragraph{Overview of the Integrated Knowledge Graph}
\label{sec:knowledge_graph}

We constructed an integrated knowledge graph to analyse the representation and interoperability of regulatory substances across datasets.
Figure~\ref{fig:knowledge_graph} provides an overview of the resulting graph structure and its main semantic layers.
The graph is organised into four layers: an input layer of 18 regulatory sources; a core layer of $\sim$17,500 regulatory substance entities represented as SKOS concepts; linking layers connecting substances to structure-based identifiers~(InChIKey), chemical classifications~(ChemOnt), and biological roles~(ChEBI); and derived layers comprising linkability tiers, embedding-based semantic clusters, and annotations ($\sim$46,600 links).

\begin{figure}[!ht]
    \centering
    \includemermaid[width=\linewidth]{images/sources/knowledge_graph.mmd}
    \caption{
        Overview of the integrated knowledge graph constructed from 18 regulatory datasets.
        The graph flows from the input layer of regulatory sources through the core layer of regulatory substances~(SKOS concepts) to linking layers~(InChIKey, ChemOnt, ChEBI) and derived layers comprising linkability tiers, semantic clusters, annotations, and an LLM-based SKOS crosswalk of 304 high-confidence triples linking non-structure substance groups to ChemOnt nodes.
    }\label{fig:knowledge_graph}
\end{figure}

\paragraph{Dataset Overview (\texorpdfstring{\hyperref[sec:contribution1]{C1}}{C1})}

Integration of 18 regulatory sources yielded 41,813 records, of which 63\% could be linked to a defined molecular structure via InChIKey.

\paragraph{Heterogeneity of Substance Definitions}
\label{sec:semantic_mismatch}

A significant share of regulatory entries (37\%) does not correspond to a defined chemical structure, consistent with Kutsarova et al.~\cite{Kutsarova2019}. These include mixtures (UVCBs), substance groups, analytical parameters, or entries without usable identifiers, reflecting differences in how “substances” are defined across regulatory contexts. 

Figure~\ref{fig:analysis4d}A shows that heterogeneity varies widely between regulatory lists: some sources consist almost entirely of structure-defined molecules, whilst others contain few structure-defined entries.

Applying the seven-tier linkability taxonomy (Section~\ref{sec:linkability}; Figure~\ref{fig:analysis4d}B) shows that 62.5\% of entries (20,909) are directly linkable via InChIKey (tier 1). Only a small fraction consists of partially resolvable groups (tiers 2--3; ${<}0.2$\%).

The remaining 37.5\% cannot be directly linked to a structure: 11.6\% have unresolved CAS numbers (tier 4), 10.8\% are mixtures or UVCBs (tier 5), ${<}0.1$\% are administrative entries (tier 6), and 14.9\% lack any usable identifier (tier 7).

More than one-third of regulatory entries cannot be mapped to a unique chemical structure, and even CAS numbers do not guarantee unambiguous identification. Consequently, structure-based integration is limited to part of the data, and the non-structure fraction requires different analytical treatment.

\begin{figure}[!ht]
    \centering
    \includegraphics[width=\linewidth]{Analysis_4bd_v_Entity_type_and_linkability_vertical.pdf}
    \caption{
        (A) Entity type composition across 18 regulatory sources; the mix of seven mutually exclusive types (Molecule, Substance group, Analytical parameter, Mixture, Regulatory entry, CAS without structure, Unclassified) differs substantially between regulatory instruments. (B) Linkability taxonomy applied to all 33,452 unique regulatory substance entries: 20,909 entries (62.5\%) are directly linked via InChIKey (tier~1); 4,981 (14.9\%) lack any usable identifier. Together, the panels show that structural linkability is unevenly distributed and that more than one-third of all entries cannot be resolved to a unique chemical structure.
    }
    \label{fig:analysis4d}
\end{figure}

\paragraph{Limitations of CAS Numbers (\texorpdfstring{\hyperref[sec:contribution2]{C2}}{C2})}

CAS numbers are not reliable unique identifiers. Bidirectional mapping analysis reveals that 195 InChIKeys are associated with more than one CAS number (up to 7 per structure), while 4 CAS numbers resolve to two distinct InChIKeys. These inconsistencies arise from differences in stereochemistry, salt forms, and context-dependent registration practices~\cite{Akhondi2012,Williams2011Quality}. CAS-based cross-dataset linking therefore introduces ambiguity rather than resolving it, making CAS numbers insufficient as a backbone for interoperability.

\paragraph{Limited Interoperability Across Regulatory Lists}

Cross-list analysis of the 17,523 structure-resolved InChIKeys shows that 70.0\% (12,270 entries) appear in only a single regulatory source, while just 30.0\% (5,253 entries) are shared across two or more lists. This limited overlap reflects the domain-specific design of regulatory instruments rather than a shared chemical reference framework.

\paragraph{Failure of Semantic Approaches for Non-Structure Substances (\texorpdfstring{\hyperref[sec:contribution3]{C3}}{C3})}
\looseness=-1 Embedding-based clustering of non-structure substance names yields semantically coherent groups but does not resolve chemical identity. While clusters reflect lexical or contextual similarity among substance descriptions, they do not correspond to well-defined molecular entities (Figure~\ref{fig:analysis89}A), confirming that neither lexical context nor pre-trained language models can recover chemical structure from regulatory substance names.

\begin{figure}[!ht]
    \centering
    \includegraphics[width=\linewidth]{Analysis_89_UMAP_and_ChemOnt_quality_ellipse.pdf}
    \caption{(A) UMAP projection of sentence embeddings for 12,499 non-structure regulatory substance entries, grouped by k-means cluster (k = 6); spatial proximity reflects semantic similarity in the 384-dimensional embedding space.  Clusters: Cl1 = Reaction masses; Cl2 = Complex substances (UVCB with structure); Cl3 = Petroleum \& coal tar fractions; Cl4 = Inorganic compounds \& metal salts; Cl5 = Trade names, codes \& biological materials; Cl6 = Polymers, fatty acids \& surfactants. (B) Distribution of top-1 cosine similarity scores between substance name embeddings and ChemOnt class label embeddings for \path{all-MiniLM-L6-v2}; the dashed line indicates the threshold ($s_1 > 0.80$) used to bound manual evaluation (132 entries, 1.1\%).}
    \label{fig:analysis89}
\end{figure}

To evaluate whether semantic similarity can recover chemical meaning, cosine similarity between sentence embeddings of substance names and ChemOnt class labels was computed for four models (\path{all-MiniLM-L6-v2}, \path{all-mpnet-base-v2}, \path{allenai/scibert\_scivocab\_uncased}, \path{dmis-lab/biobert-v1.1}). Score distributions are heavily skewed towards lower values and are not comparable across models, as embedding geometry differs between architectures (Figure~\ref{fig:analysis89}B); a fixed threshold is therefore model-dependent and not a principled quality criterion. Manual inspection of the top-200 matches per model confirmed that the large majority of high-scoring assignments are chemically incorrect; correct cases were limited to surface-form overlap between substance name and class label (e.g.\ \textit{Phenols} $\to$ \textit{Phenols}). Score magnitude is not a reliable indicator of match correctness: models with the highest score floors (SciBERT, BioBERT) performed worst on both manual inspection and quantitative validation.

To validate this failure quantitatively, the four models were evaluated against a structure-derived ChemOnt ground truth comprising 28,204 ClassyFire-assigned substances. These results, shown in Table~\ref{tab:chemont_validation}, indicate that \path{all-MiniLM-L6-v2} achieved the highest Hit@1 of 5.75\% and MRR of 0.093 — even this best result implies the correct class is retrieved in fewer than 1 in 17 attempts. All other models scored below Hit@1 = 1\%. Retrieval performance decreased further with increasing taxonomic abstraction: Hit@1 fell to 3.25\% at the subclass level and below 1\% for class, superclass, and kingdom. A hierarchy-aware weighted MRR of 0.099 was only marginally higher than the strict MRR, confirming that partial matches at higher levels of the ChemOnt hierarchy do not rescue overall classification quality.

\begin{table}[!ht]
    \centering
    \caption{Performance of embedding models for ChemOnt classification.}
    \label{tab:chemont_validation}
    \begin{tabular}{lccccc}
        \hline
        Model           & Hit@1  & Hit@3  & Hit@5  & Hit@10 & MRR   \\
        \hline
        all-MiniLM-L6-v2  & 0.0575 & 0.0986 & 0.1234 & 0.1621 & 0.093 \\
        all-mpnet-base-v2 & 0.0093 & 0.0152 & 0.0184 & 0.0247 & 0.015 \\
        SciBERT           & 0.0040 & 0.0071 & 0.0102 & 0.0131 & 0.008 \\
        BioBERT           & 0.0047 & 0.0077 & 0.0099 & 0.0150 & 0.009 \\
        Claude Sonnet~4.6$^\dagger$ & 0.392 & 0.412 & 0.414 & 0.414$^*$ & 0.402 \\
        \multicolumn{6}{l}{\footnotesize $^*$Hit@k saturates at Hit@5 because the LLM ranks over 5 hierarchy levels.} \\
        \multicolumn{6}{l}{\footnotesize $^\dagger$LLM generates a label per hierarchy level, not a ranking over the full class space; metrics} \\
        \multicolumn{6}{l}{\footnotesize \phantom{$^\dagger$}reflect per-level accuracy and are not directly comparable to embedding retrieval scores.} \\
        \hline
    \end{tabular}
\end{table}



Large language model classification substantially outperforms embedding-based retrieval: Claude Sonnet~4.6 achieves a Hit@1 of 39.2\% and MRR of 0.40, a seven-fold improvement over the best embedding model. Performance is monotonically higher at more general levels (96.5\% at kingdom, 52.3\% at superclass), and LLM-assigned confidence is a validated quality signal: high-confidence assignments reach 53.9\% accuracy, while low-confidence cases score 0\%. Nevertheless, even at its strongest, the LLM correctly classifies fewer than 2 in 5 substances at the most specific level, and application to non-structure substances yields substantially higher null rates and lower confidence, confirming that linguistic structure alone cannot substitute for chemical structure information.

\paragraph{Scale of the Interoperability Problem (\texorpdfstring{\hyperref[sec:contribution4]{C4}}{C4})}
For the more than 11,000 non-structure entries, full manual reconciliation would require on the order of hundreds of millions of pairwise comparisons. Under conservative assumptions, this corresponds to several thousand person-years of expert effort.

\begin{figure}[!ht]
    \centering
    \includegraphics[width=\linewidth]{Analysis_10c_workload_curve_per_source.pdf}
    \caption{
        Estimated pairwise curation workload for non-structure entries across the combined European regulatory dataset, broken down by source.
        Reach registrations correspond to approximately 360,000 person-days (1,800 person-years) under an assumed throughput of 100 pair assessments per person-day.
        The quadratic growth of the workload illustrates that manual reconciliation is not scalable, reinforcing the need for structure-based and ontology-driven representations at the source.
    }
    \label{fig:analysis10c}
\end{figure}



Figure~\ref{fig:analysis10c} illustrates this quadratic growth across all regulatory sources; the workload is dominated by REACH registrations (~1,800 person-years alone), and each additional non-structure entry is disproportionately costly.

\paragraph{Knowledge Graph Integration and Priority Scoring (\texorpdfstring{\hyperref[sec:contribution5]{C5}}{C5})}
Structure-defined substances can be successfully integrated across regulatory lists, chemical taxonomies (ChemOnt), and biological annotations (ChEBI) using InChIKey as a linking key. 

\begin{figure}[!ht]
	\centering
	\includegraphics[width=\linewidth]{Analysis_14e_Bubble_priority.pdf}
	\caption{Bubble chart of composite priority space: regulatory list score versus biological hazard score for all scored substances. Bubble size is proportional to the number of regulatory source memberships. Substances in the upper-right quadrant combine broad regulatory coverage with documented biological hazard evidence.}
	\label{fig:analysis14e}
\end{figure}
This enables cross-domain queries that are infeasible on fragmented regulatory datasets: for example, identifying substances that appear on three or more ECHA lists and carry a documented biological role in ChEBI, or retrieving all members of a given ChemOnt class that are subject to authorisation requirements. As an illustration of this analytical potential, a composite priority score combining regulatory list membership with ChEBI biological role annotations was computed following the ToxPi framework~\cite{Reif2010}. Figure~\ref{fig:analysis14e} plots regulatory versus biological scores for all scored substances; 2,2-bis(4-hydroxyphenyl)propane and dibutyl phthalate exemplify substances with both broad regulatory presence and multiple documented biological roles. This score is illustrative and does not constitute a formal risk assessment. Extending the knowledge graph to non-structure substance groups, LLM-based equivalence assessment identifies 31.7\% of sampled non-structure entries as ChemOnt-equivalent group entries---far above the 0.8\% detected by regular-expression pattern matching. Of these, 304 high-confidence SKOS mapping triples were validated (285~\texttt{skos:broadMatch}, 18~\texttt{skos:exactMatch}, 1~\texttt{skos:narrowMatch}), forming a directly reusable legislative--ChemOnt crosswalk. The dominance of \texttt{skos:broadMatch} (92.8\%) reflects that regulatory group entries typically define constrained subsets of ChemOnt classes (e.g.\ \textit{fatty acids, C16--22} $\subset$ \textit{Fatty acids and conjugates}), making their scope explicit and their member substances enumerable via the ChemOnt hierarchy. The remaining 68\% of non-structure entries could not be reliably mapped, underscoring the limits of post-hoc taxonomic alignment.





